% !TEX root = BUT-TD1-R3D16-GET-PRICES.tex
\input{../inc/encommun.tex}
\graphicspath{{images/}{./images/}}

\usepackage{minted}
\PassOptionsToPackage{hyphens}{url}\usepackage{hyperref}

\title{BUT TD1 R3D16\\ 
(Automatisation de la récupération des prix d'AWS)}
\date{Octobre 2025} % TD prévu prévu le ?/10 2014
\author{Jean-Marc Pouchoulon}
% pour afficher ( ou pas) les réponses dans le document
%\printanswers
\noprintanswers

\begin{document}-è
\maketitle
Ce TD a pour but de prendre connaissance du modèle de prix des instances "COMPUTE" de aws.
Le calcul du prix d'une instance de machine virtuelle dans le Cloud est une demande récurrente afin de comparer les prix des instances 
en fonction des régions AWS et du type d'instance.  


%%%%%%%%%%%%%%%%%%%%%%%%%%%%%%%%%%%%%%%%%%%%%%%%%%%%%%%%%%%%%%%%
\section{Obtention des prix au format json des instances "COMPUTE" d'aws via curl}

\subsection{Installation d'un container de prix}


Installez via docker-compose le container "EC2 shop"  permettant de récupérer des prix.
\\
voir: \url{https://github.com/yeo/ec2.shop} et \url{https://ec2.shop/}

\begin{questions}
\question Récupérez les prix des instances m5 et m6 de la region us-east puis eu-west-1 et sauvegardez les dans un fichier tiny.json.
\question  Récupérez les prix des instances t1 et t2.micro de la region eu-west-1 sauvegardez les dans un fichier medium.json.
\question à quoi correspondent les données extraites ?
\question Quel est la différence de prix au mois entre une m6a.48xlarge et une t1.micro ?
\begin{solution}
\begin{bashcode}
curl -H 'accept: json' 'http://localhost:6001?region=us-east-1&filter=t1,t2&sort=vCPUs'|jq '.' > tiny.json
curl -H 'accept: json' 'http://localhost:6001?region=us-east-1&filter=m5,m6&sort=vCPUs'|jq '.'> medium.json
jq -r '.Prices[]|select(.InstanceType=="t1.micro")|.InstanceType,.MonthlyPrice' tiny.json
t1.micro
14.60
> jq -r '.Prices[]|select(.InstanceType=="m6a.48xlarge")|.InstanceType,.MonthlyPrice' medium.json
m6a.48xlarge
6054.911999999999
nu
# nushell
open tiny.json|flatten|flatten|sort-by VCPUS
open tiny.json|flatten|flatten|sort-by Memory 
            | insert MemStr {$in.Memory |split column ' '}
            | flatten --all|flatten --all 
            | reject column2 Memory
            | insert MemNumGiB {$in.column1|into decimal} 
            | reject column1
            | sort-by MemNumGiB
            | move MemNumGiB --before Storage

\end{bashcode}
\end{solution}

\end{questions}



\section{Transformation des fichiers json en csv }

Utilisez l'utilitaire jq pour répondre aux questions.
Pour afficher
\begin{bashcode}
jq -r '.' medium.json
jq -r '.' tiny.json
\end{bashcode}


\begin{questions}
       \question à quoi sert l'option "-r" ?
       \question Afficher juste le contenu de 'Prices' en json' ?
       \question quelles sont les clefs des entrées (utilisez '|keys|@csv') ?
       \question n'affichez qu'une seule ligne pour les clefs qui constituera les en-têtes de votre fichier csv.
       \question utilisez l'ordre "join" pour produire les lignes du fichier csv.
       \question A l'aide de bash assemblez le header et les lignes pour créer le csv. Vérifiez que le fichier csv produit est importable dans libre office.
       \question trouvez une autre solution pour extraire.
% \begin{solution}
% \begin{yamlcode}
%        https://qmacro.org/blog/posts/2022/05/19/json-object-values-into-csv-with-jq/

%        ## curl des prix region us-east-1
%        ```
%        curl -H 'accept: json' 'http://localhost:6001?region=us-east-1&filter=m5,m6&sort=vCPUs'|jq '.' m5.json
%        ```
       
       
%        ## Pour les clefs du csv
       
%        ```
%        > jq -r '.Prices[]| [.InstanceType, .Memory, .VCPUS, .Storage, .Network, .Price, .MonthlyPrice, .SpotPrice]\|@csv' m5.json   
%        ```
       
       
%        ```
%        jq -r '.Prices[]| keys|@csv' medium.json
%        "Cost","InstanceType","Memory","MonthlyPrice","Network","SpotPrice","Storage","VCPUS"
%        "Cost","InstanceType","Memory","MonthlyPrice","Network","SpotPrice","Storage","VCPUS"
%        "Cost","InstanceType","Memory","MonthlyPrice"",","Network","SpotPrice","Storage","VCPUS"
%        "Cost","InstanceType","Memory","MonthlyPrice","Network","SpotPrice","Storage","VCPUS"
%        ```
       
%        ## juste les champs . On sélectionne la première ligne, on récupère les clefs dans $k et on affiche $k
       
%        ```
%         jq -r '.Prices[0]| keys' medium.json
%        [
%          "Cost",
%          "InstanceType",
%          "Memory",
%          "MonthlyPrice",
%          "Network",
%          "SpotPrice",
%          "Storage",
%          "VCPUS"
%        ]
%        ❯ jq -r '.Prices[0]| keys[]' m5.json
%        Cost
%        InstanceType
%        Memory
%        MonthlyPrice
%        Network
%        SpotPrice
%        Storage
%        VCPUS
       
%        > jq -r '.Prices[0]| keys[] as $k|$k' m5.json
%        Cost
%        InstanceType
%        Memory
%        MonthlyPrice
%        Network
%        SpotPrice
%        Storage
%        VCPUS
       
%        >  jq -r '.Prices[0]| keys|@csv' m5.json
%        "Cost","InstanceType","Memory","MonthlyPrice","Network","SpotPrice","Storage","VCPUS"
       
       
%        ```
%        ## Pour les valeurs du csv
       
%        ### un moyen simple
       
%        ```
%        jq -r '.Prices[]| join(",")' medium.json 
       
%        m5ad.large,8 GiB,2,1 x 75 NVMe SSD,Up to 10 Gigabit,0.103,75.19,0.0344
%        m5n.4xlarge,64 GiB,16,EBS only,Up to 25 Gigabit,0.952,694.9599999999999,0.4110
%        m6g.2xlarge,32 GiB,8,EBS only,Up to 10 Gigabit,0.308,224.84,0.1696
%        m6i.metal,512 GiB,128,EBS only,50000 Megabit,6.144,4485.12,2.2852
%        m6id.24xlarge,384 GiB,96,4 x 1425 SSD,37500 Megabit,5.6952,4157.496,1.7139
%        m5dn.2xlarge,32 GiB,8,1 x 300 NVMe SSD,Up to 25 
%        ...
%        ### moyen plus complexe
       
%        #### Récupération des valeurs
       
%        ```
%        jq -r '.Prices[]| [.InstanceType, .Memory, .VCPUS, .Storage, .Network, .Price, .MonthlyPrice, .SpotPrice]' m5.json|jq
       
       
%        [
%          [
%            "m6i.2xlarge",
%            "32 GiB",
%            8,
%            "EBS only",
%            "Up to 12500 Megabit",
%            null,
%            280.32,
%            "0.1759"
%          ],
%          [
%            "m5a.8xlarge",
%            "128 GiB",
%            32,
%            "EBS only",
%            "Up to 10 Gigabit",
%            null,
%            1004.4799999999999,
%            "0.5915"
%        ...
       
%        On peut aussi utiliser map
%        ```
%        jq -r '.Prices| map([.InstanceType, .Memory, .VCPUS, .Storage, .Network, .Price, .MonthlyPrice, .SpotPrice])' m5.json
%        ```
       
       
%        ##### passage au csv
%        ```
%        jq -r '.Prices[]| [.InstanceType, .Memory, .VCPUS, .Storage, .Network, .Price, .MonthlyPrice, .SpotPrice]|@csv' m5.json
%        ```
       
%        ```
%        jq -r '.Prices| map([.InstanceType, .Memory, .VCPUS, .Storage, .Network, .Price, .MonthlyPrice, .SpotPrice])|.[]|@csv' m5.json
%        ```
       
       
%        ### formes générique
       
%        ```
%        jq -r '.Prices[] | [keys[] as $k | .[$k]]|@csv' m5.json 
       
%        0.384,"m6i.2xlarge","32 GiB",280.32,"Up to 12500 Megabit","0.1759","EBS only",8
%        1.376,"m5a.8xlarge","128 GiB",1004.4799999999999,"Up to 10 Gigabit","0.5915","EBS only",32
%        6.528,"m5dn.24xlarge","384 GiB",4765.44,"100 Gigabit","1.6941","4 x 900 NVMe SSD",96
%        0.904,"m5d.4xlarge","64 GiB",659.9200000000001,"Up to 10 Gigabit","0.3574","2 x 300 NVMe SSD",16
%        0.077,"m6g.large","8 GiB",56.21,"Up to 10 Gigabit","0.0464","EBS only",2
%        0.1808,"m6gd.xlarge","16 GiB",131.98399999999998,"Up to 10 Gigabit","0.0714","1 x 237 NVMe SSD",4
%        1.648,"m5ad.8xlarge","128 GiB",1203.04,"Up to 10 Gigabit","0.7226","2 x 600 NVMe SSD",32
%        0.226,"m5d.xlarge","16 GiB",164.98000000000002,"Up to 10 Gigabit","0.0776","1 x 150 NVMe SSD",4
%        0.192,"m5.xlarge","16 GiB",140.16,"Up to 10 Gigabit","0.0686","EBS only",4
%        2.1696,"m6gd.12xlarge","192 GiB",1583.808,"20 Gigabit","0.8570","2 x 1425 NVMe SSD",48
%        7.5936,"m6id.32xlarge","512 GiB",5543.328,"50000 Megabit","2.2852","4 x 1900 SSD",128
%        0.113,"m5d.large","8 GiB",82.49000000000001,"Up to 10 Gigabit","0.0360","1 x 75 NVMe SSD",2
%        5.712,"m5n.metal","384 GiB",4169.76,"100 Gigabit","1.6323","EBS only",96
%        0.616,"m6g.4xlarge","64 GiB",449.68,"Up to 10 Gigabit","0.3584","EBS only",16
%        0.096,"m6i.large","8 GiB",70.08,"Up to 12500 Megabit","0.0366","EBS only",2
%        0.119,"m5n.large","8 GiB",86.86999999999999,"Up to 25 Gigabit","0.0340","EBS only",2
%        ...
       
%        ```
%        Redirection vers le csv
%        ```
%        jq -r '.Prices[] | [keys[] as $k | .[$k]]|@csv' m5.json|less>> m5.csv
%        ```
% \end{yamlcode}

% \end{solution}  
\end{questions}

\end{document}